\documentclass{article}
\usepackage{graphicx}

\usepackage[utf8]{inputenc}     % pdfLaTeX
\usepackage[T1]{fontenc}
\usepackage[magyar]{babel}

\title{Hallás}
\author{Lengyel Zoltán Péter, Rovenszky Robin}
\date{Legutóbb frissítve: 2026. Március 25.}

\begin{document}

\maketitle
\tableofcontents
\newpage

\section{Fül}
% insert fül figures here
\begin{itemize}
    \item A hallószerv három részre osztható: 
    \begin{itemize}
        \item Külső fül \begin{itemize}
            \item Fülkagyló
            \item Külső hallójárat
        \end{itemize}
        \item Középfül \begin{itemize}
            \item Dobhártya 
            \item 3 hallócsont \begin{itemize}
                \item Kalapács (dobhártyához kapcsolódik)
                \item Üllő 
                \item Kengyel
            \end{itemize}
        \end{itemize}
        \item Belső fül \begin{itemize}
            \item Csiga - folyadékkal tültütt üregrendszer \begin{itemize}
                \item Ovális ablak
                \item Felső járat
                \item Középső járat -  belsejében mechanoreceptorok (Corti-féle szerv: érzékszőrök + fedő)
                \item Alsó járat
                \item A hang elindul a felső járaton, a csiga végén visszafordul és az alsó járaton jön vissza
                \item Kerek ablak
            \end{itemize}
            \item Fülkürt (garattal van összeköttetésben, dobhártya két oldala közti légnyomást egyenlíti ki - nyelés)
        \end{itemize}
    \end{itemize}

    \item 400Hz-22000Hz között érzékel \begin{itemize}
        \item Mély hangok - a csiga csúcsa környékén váltanak ki ingerületet
        \item Magas hangok -  ovális ablakhoz közel váltanak ki ingerületet
    \end{itemize}
\end{itemize}
%dB table here

\subsection{Egyensúly érzékelő szervrendszer}
\begin{itemize}
    \item Mechanoreceptorok \begin{itemize}
        \item Ampulla - ívjárat alapi része közelében \begin{itemize}
            \item Szőrsejtek 
            \item Kupula - szőrsejteket körbevevő kocsonyás kúp
        \end{itemize}
    \end{itemize}
    \item Tömlőcske és zsákocska között \begin{itemize}
        \item Nagyobb, szélesebb alapon elhelyezkedő sejtek, tetején mészkristályok nyomják a szőrsejteket
    \end{itemize}
\end{itemize}

\end{document}